\documentclass{article}

\usepackage{amsmath,amssymb}
\usepackage{xurl}
\usepackage[hidelinks]{hyperref}
\setlength{\emergencystretch}{2em}

% Shared commands for notes in docs/paper-gaps/.

\DeclareUnicodeCharacter{2080}{\ifmmode{}_{0}\else\textsubscript{0}\fi}
\DeclareUnicodeCharacter{2081}{\ifmmode{}_{1}\else\textsubscript{1}\fi}
\DeclareUnicodeCharacter{2082}{\ifmmode{}_{2}\else\textsubscript{2}\fi}
\DeclareUnicodeCharacter{2099}{\ifmmode{}_{n}\else\textsubscript{n}\fi}

\newcommand{\Lean}{\textsc{Lean}}
\ExplSyntaxOn
\NewDocumentCommand{\leanid}{m}{
  \ifmmode
    \textsf{\detokenize{#1}}
  \else
    \group_begin:
    \urlstyle{sf}
    \tl_set:Nn \l_tmpa_tl {#1}
    \tl_replace_all:Nnn \l_tmpa_tl {\_} {_}
    \exp_args:NV \path \l_tmpa_tl
    \group_end:
  \fi
}
\ExplSyntaxOff
\providecommand{\tr}{\operatorname{tr}}
\newcommand{\tnleanIssueBase}{https://github.com/LionSR/TNLean/issues/}
\newcommand{\tnleanPrBase}{https://github.com/LionSR/TNLean/pull/}
\newcommand{\tnleanDocBase}{https://sirui-lu.com/TNLean/docs/}
\newcommand{\ghissue}[1]{\href{\tnleanIssueBase#1}{GitHub issue \##1}}
\newcommand{\ghpr}[1]{\href{\tnleanPrBase#1}{PR \##1}}
\newcommand{\doclink}[1]{\href{\tnleanDocBase#1.html}{\path{#1}}}

% Dirac notation and the identity operator, as in blueprint/src/macros/common.tex.
% \providecommand keeps note-local definitions authoritative.
\usepackage{dsfont}
\providecommand{\ket}[1]{|#1\rangle}
\providecommand{\bra}[1]{\langle #1|}
\providecommand{\braket}[2]{\langle #1 | #2 \rangle}
\providecommand{\ketbra}[2]{|#1\rangle\!\langle #2|}
\providecommand{\Id}{\mathds{1}}
\providecommand{\one}{\mathds{1}}
\providecommand{\C}{\mathbb{C}}
\providecommand{\MN}[1]{M_{#1}(\C)}
\providecommand{\MD}{M_D(\C)}

% External referencing for paper sources.  The build helper writes sanitized
% auxiliary files under build/paper-aux/ and excludes duplicated source labels.
\usepackage{xr}

% Import a generated paper auxiliary file with a caller-supplied label prefix.
% The candidate paths support builds from docs/paper-gaps/, the repository root,
% and build/paper-aux/ itself.
\newcommand{\importpaperaux}[2]{%
  \IfFileExists{../../build/paper-aux/#2.aux}{%
    \externaldocument[#1]{../../build/paper-aux/#2}%
  }{%
    \IfFileExists{build/paper-aux/#2.aux}{%
      \externaldocument[#1]{build/paper-aux/#2}%
    }{%
      \IfFileExists{#2.aux}{%
        \externaldocument[#1]{#2}%
      }{}%
    }%
  }%
}

% General external reference macros
\newcommand{\paperref}[2]{\ref{#1-#2}}
\newcommand{\papereqref}[2]{(\ref{#1-#2})}

% CPSV16 source (1606.00608 / MPDO-22-12-17-2.tex) external document and shortcuts
\importpaperaux{cpsv16-}{1606.00608/MPDO-22-12-17-2-tnlean}
\newcommand{\cpsvref}[1]{\paperref{cpsv16}{#1}}
\newcommand{\cpsveqref}[1]{\papereqref{cpsv16}{#1}}

% Verdict marker: \gapnote{<kind>}{<status>}. Kind is the policy
% classification (clarification, local-correction, scope-restriction,
% unfaithful, false-source, open-gap); status is open, wip (elimination
% underway), resolved, or historical. Machine-read by the site index;
% prints a verdict line.
\newcommand{\gapnote}[2]{\par\noindent\textbf{Verdict:} #1 (#2).\par}


\title{Scope of the First Choi-Type Positivity Subcase (Wolf Example~3.1)}
\author{}
\date{2026-07-02}

\begin{document}
\maketitle

\section{Source statement}

Wolf's Chapter~3, Example~3.1 introduces the Choi-type maps
\cite{Wolf2012Quantum}\footnote{Local source:
\path{Notes/WolfNoteTexSource/ch03_positive_not_completely.tex}, lines
357--365.}  With \(D\) the diagonal projection on \(\mathcal M_d\) and
\(U_{k0}\) the cyclic shift, the map is
\begin{align}
  T_C(X)
  =(d-n)D(X)-X+\sum_{k=1}^{n}D(U_{k0}XU_{k0}^{\dagger}).
  \tag{3.20}
\end{align}
The source states that \(T_C\) is positive and indecomposable for every
\(1\le n\le d-2\).  No nonzero-coordinate hypothesis is present.

\section{Restricted statement now proved}

The present development proves only the first rank-one subcase.  For \(d=3\),
\(n=1\), and an arbitrary vector \(v=(v_0,v_1,v_2)\), it proves
\begin{align}
  T_C(|v\rangle\langle v|)\ge 0.
\end{align}
The scalar input is the boundary-inclusive cyclic reciprocal inequality
\begin{align}
  \frac{x}{2x+z}+\frac{y}{2y+x}+\frac{z}{2z+y}\le 1,
  \qquad x,y,z\ge 0,
\end{align}
with zero-denominator summands read as \(0\).  Away from the boundary this is
proved by clearing denominators and applying AM--GM to \(x^2y,y^2z,z^2x\); the
boundary cases reduce to a single half-bound.\footnote{Formal declarations:
\leanid{Matrix.choiType\_cyclic\_reciprocal\_three\_one\_of\_nonneg} and
\leanid{Matrix.choiTypeMap\_vecMulVec\_posSemidef\_three\_one}
in \doclink{TNLean/Channel/ChoiTypeMap}.}

\section{Missing hypotheses and elimination plan}

This is a scope restriction, not the source theorem.  The full positivity
statement still requires the nonnegative cyclic reciprocal estimate
\begin{align}
  \sum_{i\in\mathbb Z/d\mathbb Z}
  \frac{x_i}{(d-n)x_i+\sum_{k=1}^{n}x_{i-k}}\le 1
  \qquad (x_i\ge0,\;1\le n\le d-2),
\end{align}
with the usual zero-term convention.  Once this estimate is available, it
should be applied to \(x_i=|v_i|^2\) to obtain the rank-one statement for the
full parameter range.  A positive-semidefinite decomposition will then extend
the rank-one result to positivity of \(T_C\) on all positive semidefinite
matrices.  This is tracked by \ghissue{3622}.  The indecomposability assertion
of the same source paragraph is separate and is tracked by \ghissue{3623}.

\section{Classification}

\emph{Scope restriction.}  The current theorem is mathematically correct as a
\(3\times3\) rank-one subcase, but it should not be cited as the formalization
of Wolf's full Choi-type positivity theorem.

\bibliographystyle{alpha}
\bibliography{references}

\end{document}
